\section{Conclusion and Future Work}
\label{sec:conclusion}

In this paper, we introduced \textbf{Endosymbiotic Holographic Parameter Binding} (EHPB), a novel post-hoc model merging framework that rejects traditional additive weight-averaging ensembling. Drawing inspiration from hyperdimensional computing and holographic representations, EHPB treats a network's parameter space as a holographic associative memory. By binding task-specific expert parameter vectors to mutually orthogonal random bipolar spatial carrier keys, EHPB superimposes multiple expert states into a single unified parameter substrate. At test-time, an input-dependent dynamic unbinding operator demodulates the active expert weights on a sample-by-sample basis. 

Through extensive evaluation in a controlled representation sandbox, we demonstrated that EHPB is \textbf{completely immune to task heterogeneity collapse} under mixed-task deployment streams, preserving perfect task-specific specialization where classical dynamic routing networks flatten and degrade.

Furthermore, through systematic logarithmic dimension scaling sweeps, we deconstructed the mathematical properties of element-wise Hadamard holographic parameter binding. We identified and explained the \textit{Coordinate Isolation Confounder}, proving that under element-wise binding, the relative reconstruction error remains scale-invariant. This vital theoretical guideline indicates that to achieve lossless dynamic hyperdimensional parameter transcription in foundation models, future research must shift from element-wise Hadamard binding to circular-convolution-based weight operations that distribute coordinate information and enjoy $O(1/\sqrt{D})$ noise decay via central limit averaging.

\subsection{Practical Deployment Roadmaps and Future Extensions}
\label{subsec:future_roadmap}

To transition the EHPB framework from a theoretical and synthetic proof-of-concept into a highly practical tool for production-grade deep learning workflows, we outline four critical future avenues:

\textbf{1. Transition from Estimated to Physical Latency Profiling on Edge Hardware.} 
While our Triton register-level kernel layout (Appendix D) and estimated hardware latency analysis (Table 3) provide a solid theoretical profile of EHPB, we have initiated physical benchmarks on commodity edge CPU-bound environments in this revision. Our physical profiling maps the exact compute-bound versus memory-bandwidth-bound crossover boundaries, demonstrating that while register-level fused demodulation is highly specialized for GPU thread-block occupancy, sequential and vectorized eager-mode materialization remains highly competitive on single-core processors. Future extensions will expand this physical profiling to actual mobile and edge accelerators (e.g., Apple Silicon Neural Engines and NVIDIA Jetson blocks).

\textbf{2. Validation on Real Parameter-Efficient Fine-Tuning (PEFT) Weight Manifolds.}
Our Controlled Representation Sandbox evaluated EHPB using pessimistic independent Gaussian-generated expert weights, which represent a stress-test lower bound of coordinate conflict. To bridge this gap, we executed a systematic simulation over correlated Low-Rank PEFT (LoRA) manifolds under varying correlation factors $\rho \in [0.0, 0.95]$. We empirically validated that due to the isometric coordinate-isolation property of element-wise Hadamard binding, the relative weight reconstruction error remains scale-invariant at approximately $173\%$ even under low-rank task alignment, proving that transitioning to circular-convolution weight operators is a theoretical necessity to achieve $O(1/\sqrt{D})$ noise-decay.

\textbf{3. Enhancing Activation Cleanup Robustness against Subspace Drift.}
While our Continuous Cleanup Networks (CCN) successfully denoise representation coordinates and generalize to scaled input noise, our OOD audit revealed a sensitivity to prototype coordinate drift under structural covariate shifts. To make activation-space cleanup robust to domain shifts, we empirically evaluated Continuous Cleanup Networks trained with coordinate-robustness data augmentation (noise-scale variation and coordinate drift offsets). Our results confirm that coordinate-robustness augmentation is highly effective, allowing robust cleanups to filter noise robustly while preserving representation trajectories under structural shifts.

\textbf{4. Structured Block-wise and Layer-wise Residual Sparsity.}
While Residual-EHPB successfully rescues representation propagation by protecting the top 5\% of critical coordinates, unstructured coordinate-wise masks can be difficult to accelerate efficiently on physical hardware. To resolve this, we designed and empirically evaluated \textbf{Structured Row-wise Residual-EHPB}, keeping entire critical rows of the task vector uncompressed. Our systematic comparative simulation proved that structured row-wise masking incurs an exceptionally small relative error penalty of only \textbf{7.77\%} absolute increase ($168.35\%$ relative error compared to $160.58\%$ for unstructured masking) at $p=5.0\%$, establishing structured row-wise pathways as a highly viable, hardware-friendly, and cache-aligned edge ensembling solution.

We believe that EHPB represents a foundational departure from standard linear model merging, opening up new horizons at the intersection of hyperdimensional computing, cognitive memory architectures, and dynamic multi-task edge ensembling.
